\originalchapter{级数、闭环与对偶}
\label{ch:duality}
\source{8.334 L15--L17。}
\originalsection{高温展开与低温畴壁}
Ising键$\e^{Ks_is_j}=\cosh K(1+s_is_j\tanh K)$，逐键相乘后对所有自旋求和，任何顶点若有奇数入边则该自旋求和为0。故$Z=2^N(\cosh K)^{N_b}\sum_{\rm even\ graphs}(\tanh K)^{|G|}$，图在每顶点偶度，在二维可分成闭环但交叉与重叠的计数须保留。

低温从全同自旋出发，翻一个畴的能量只随边界长度：每条不满足键代价$2J$。于是低温图权重$\e^{-2K\ell}$，图画在对偶格。Peierls估计在二维把穿过某邻域的长度$\ell$闭环数上界为常数乘$3^\ell$，若$3\e^{-2K}<1$，长畴壁总概率小，可证明足够低温有自发磁化。这个粗界证明存在有序区，不给精确$T_c$。
\originalsection{一维传递矩阵}
含场$h$时$T_{ss'}=\exp[Kss'+h(s+s')/2]$，$Z=\Tr T^N$。本征值$\lambda_\pm=\e^K\cosh h\pm\sqrt{\e^{2K}\sinh^2h+\e^{-2K}}$，无限链$f=-\kb T\log\lambda_+$。有限$K$最大根光滑且唯一，因此无有限温非解析相变。零场$\lambda_-/\lambda_+=\tanh K$，再次得到相关长度。

高温闭环或低温缺陷在一维都可完整求和，适合核对展开系数和有限尺寸实现。把一维无相变推到二维的错误在于域壁熵随长度的增长和环结构已变化。
\originalsection{二维自对偶与三维规范对偶}
方格的低温畴壁与对偶格高温偶图一一对应，令$\e^{-2K^*}=\tanh K$。扣去显式$2^N,\cosh K$等解析因子，$Z(K)$映到$Z(K^*)$。自对偶条件$\sinh2K_c=1$给$K_c=\frac12\log(1+\sqrt2)$，但只凭自对偶还须知道相变位置与唯一性；下章精确自由能核对此唯一临界奇异。

三维Ising低温图是闭合的二维畴壁面，而非二维闭线。$\mathbb Z_2$格点规范模型边变量$u_l=\pm1$、作用$-K_g\sum_p\prod_{l\in\partial p}u_l$的高温展开也给闭面，因为对每条边求和要求相邻选面数偶。因此$\e^{-2K_{\mathrm{Ising}}}=\tanh K_g$联系三维Ising与对偶规范模型。不能把三维模型错误称为与同一三维Ising自对偶。
\originalsection{随机游走生成与phantom环}
Gaussian格点场的二次矩阵写$A=I-w\mathsf J$，$\mathsf J$为邻接矩阵，则$\log Z=-\frac12\log\det A+\mathrm{const}=\frac12\sum_{l\ge1}w^l\Tr\mathsf J^l/l+\mathrm{const}$。$\Tr\mathsf J^l$计起点相同的闭游走，$1/l$除去起点重复，$1/2$与无向或实场计数相联。这些游走可自交、重复走边，故称phantom环。

它可总结高温图的Gaussian近似，解释高维临界的均场结构，但不是把每个Ising偶图精确计一次。二维精确求和还需转角相位，使交叉路径的过计数抵消（Kac--Ward/Pfaffian形式）。图的权重与符号结构不能只用“所有闭环都有同样权”替代。
